\documentclass[10pt,a4paper]{amsart}
\usepackage[margin=19mm]{geometry}
\usepackage{amsmath,amssymb,mathtools}
\usepackage[scr=rsfs,cal=euler]{mathalfa}
\usepackage{xfrac}
\usepackage[unicode,hidelinks,bookmarks=false]{hyperref}
\newcommand{\ie}{i.e.\ }
\newcommand{\cf}{cf.\ }
\newcommand{\resp}{respectively\ }
\newcommand{\Subsection}{\subsection}
\providecommand{\nobreakdash}{\nobreak\hbox{-}}
\setlength{\emergencystretch}{2em}
\multlinegap0pt
\usepackage{tikz}
\usepackage{enumerate}
\usepackage[shortlabels]{enumitem}
\usepackage{mathtools}
\usepackage{amssymb}
\usepackage{dsfont}
\usepackage{float}
\usepackage{bm}
\newtheorem{proposition}{Proposition}
\title{\bf\Large Fixed-density profiles for the semi-induced 4-vertex star}
\author{Jinghua Deng, Jianfeng Hou}
\date{Source: arXiv:2606.23351v1 (CC BY 4.0)}
\begin{document}
\maketitle

\noindent\emph{Source and licence.} \bf\Large Fixed-density profiles for the semi-induced 4-vertex star. Version arXiv:2606.23351v1 (CC BY 4.0), distributed by arXiv under the Creative Commons Attribution 4.0 International licence, http://creativecommons.org/licenses/by/4.0/, as recorded in the arXiv licence metadata for this exact version and shown on the abstract page https://arxiv.org/abs/2606.23351v1. Full licence notice: \texttt{licenses/arxiv-2606.23351-CC-BY-4.0.txt}.\par\smallskip

\noindent\textit{Editorial scope.} Counted unit: the article's proposition lem:split-quasiclique-comparison (source0.tex lines 274--277). The article's complete original proof of it is delivered with it (source0.tex lines 3067--3389, one \texttt{proof} environment of 323 lines, which the article places away from the statement). No mathematics is written, repaired or re-derived: the statement and the proof are the article's own lines, in the article's own notation, and no retained line is substituted or re-flowed.  Retained uncounted as the source-local context the proof needs: lines 269--271.  Nothing was omitted from any retained span, so the package carries no omission record.  No retained line contains a citation, so the wrapper carries no bibliography and no bibliography entry.  No retained line contains a figure environment, an image-inclusion command or drawing code, so no image is delivered and the package declares no asset dependency.  Editorial formatting only, in addition: an AMS \texttt{amsart} preamble that restates the article's own declarations and macros used by the retained lines, namely the environments \texttt{proposition} on separate counters, together with the standard packages those lines need.  The retained environments print in this document as Proposition 1 in order of appearance, whereas the article prints the same items with the numbering of its own full document.  The complete native source of the article as distributed in the arXiv e-print for this exact version is bundled unchanged as \texttt{sources/203424/source0.tex}.
\begin{align}\label{eq:switching-equation}
        6t^5-384t^4+338t^3-89t^2-4t+2=0.
\end{align}

\begin{proposition}\label{lem:split-quasiclique-comparison}
If $0<\beta<t_*^2$, then $s(\beta)>\beta^{3/2}(1-\sqrt{\beta})$; if $t_*^2\le \beta\le 1/4$, then 
$s(\beta)=\beta^{3/2}(1-\sqrt{\beta})$.
\end{proposition}

\begin{proof}[Proof of Proposition \ref{lem:split-quasiclique-comparison}]

 Let
 \[
 Q(t):=6t^5-384t^4+338t^3-89t^2-4t+2
 \]
 be the polynomial in \eqref{eq:switching-equation}. %We first prove that
% \(Q\) has a unique root in \((0,1/3)\). Clearly
% \[
% Q(0)=2>0,\qquad Q(1/3)=-\frac{115}{81}<0.
% \]
% Moreover,
% \[
% Q'(t)
% =
% -4(1-3t)^4
% -226t(1-3t)^3
% -804t^2(1-3t)^2
% -690t^3(1-3t)
% -582t^4 .
% \]
% For \(0\le t\le 1/3\), every term on the right-hand side is
% nonpositive, and at least one term is strictly negative. Hence
% \[
% Q'(\gamma)<0\qquad (0\le \gamma\le 1/3).
% \]
% Thus \(Q\) is strictly decreasing on \([0,1/3]\), and so it has exactly one
% root in \((0,1/3)\). Denote this root by \(t_*\).


Recall that $c(\beta)=\beta^{3/2}(1-\sqrt{\beta})$. It remains to compare \(s(\tau^2)\) with \(c(\tau^2)\). Put
\[
\Phi(x,y):=xy^2(1-y)+y(x+y)^2(1-x-y).
\]
Let \(0<\tau\le 1/2\). By definition,
\[
s(\tau^2)
=
\max\{\Phi(x,y): x,y\ge0,\ x+y\le1,\ 2xy+y^2=\tau^2\}.
\]
The point \((x,y)=(0,\tau)\) is feasible and gives
\[
\Phi(0,\tau)=\tau^3(1-\tau)=c(\tau^2).
\]


Let \((x,y)\) be feasible. Since \(\tau>0\), we have \(y>0\). Write
$y=\tau\lambda $. The constraint \(2xy+y^2=\tau^2\) gives
\begin{align}\label{eq:x}
x=\frac{\tau(1-\lambda^2)}{2\lambda}.
\end{align}
The feasibility conditions \(x\ge0\) and \(x+y\le1\) are equivalent to
\begin{equation}\label{eq:app-transition-feasible}
0<\lambda\le1,
\qquad
\tau\le \frac{2\lambda}{1+\lambda^2}.
\end{equation}
Conversely, every \(\lambda\) satisfying
\eqref{eq:app-transition-feasible} gives a feasible pair by the above
formulae.


Substituting~\eqref{eq:x} and $y=\tau\lambda$
into \(\Phi(x,y)\), a direct expansion gives
\begin{equation}\label{eq:app-transition-compare}
\Phi\left(\frac{\tau(1-\lambda^2)}{2\lambda},\tau\lambda\right)
-
c(\tau^2)
=
\frac{\tau^3(1-\lambda)^2}{8\lambda^2}T_\tau(\lambda),
\end{equation}
where
\[
T_\tau(\lambda)=B(\lambda)-\tau D(\lambda),
\qquad
B(\lambda)=-2\lambda((\lambda+1)^2-2),
\qquad
D(\lambda)=(1+\lambda)^2(1-3\lambda^2).
\]
For \(0<\lambda<1\), the prefactor in
\eqref{eq:app-transition-compare} is positive, so the sign of the difference
is the sign of \(T_\tau(\lambda)\). At \(\lambda=1\), the difference is zero.


On the interval \(0<\lambda<\sqrt{2}-1\), both \(B(\lambda)\) and
\(D(\lambda)\) are positive. Define
\begin{equation}\label{eq:app-transition-gdef}
g(\lambda):=\frac{B(\lambda)}{D(\lambda)}
=
\frac{-2\lambda((\lambda+1)^2-2)}
{(1+\lambda)^2(1-3\lambda^2)}.
\end{equation}
Then
\[
T_\tau(\lambda)=D(\lambda)(g(\lambda)-\tau),
\]
and therefore
\begin{equation}\label{eq:app-transition-Tsign}
T_\tau(\lambda)>0
\quad\Longleftrightarrow\quad
\tau<g(\lambda)
\qquad
(0<\lambda<\sqrt{2}-1).
\end{equation}


A direct differentiation gives
\begin{equation}\label{eq:app-transition-gderiv}
g'(\lambda)
=
-\frac{2H(\lambda)}
{(1+\lambda)^3(3\lambda^2-1)^2},
\end{equation}
where
\[
H(\lambda)=3\lambda^5+9\lambda^4-8\lambda^3+5\lambda-1.
\]
Moreover,
\[
H'(\lambda)=15\lambda^4+36\lambda^3-24\lambda^2+5.
\]
For \(0\le\lambda\le\sqrt{2}-1\), since \(\sqrt{2}-1<1/2\), we have
\[
H'(\lambda)\ge 5-24\lambda^2
\ge 5-24(\sqrt{2}-1)^2>0.
\]
Also
$H(0)=-1$ and
$H(\sqrt{2}-1)=80-56\sqrt{2}>0.$ 
Hence \(H\) has a unique zero \(\lambda_0\) in
\((0,\sqrt{2}-1)\). Since the denominator in
\eqref{eq:app-transition-gderiv} is positive, it follows that \(g\) is
increasing on \((0,\lambda_0)\) and decreasing on
\((\lambda_0,\sqrt{2}-1)\). Consequently,
\begin{equation}\label{eq:app-transition-gmax}
g(\lambda)\le g(\lambda_0)
\qquad
(0<\lambda<\sqrt{2}-1).
\end{equation}
Put
\[
m:=g(\lambda_0).
\]


We next identify \(m\). First, \(0<m<1/3\). Indeed, \(m>0\) because
\(B(\lambda)\) and \(D(\lambda)\) are positive on
\(0<\lambda<\sqrt{2}-1\). Also,
\[
D(\lambda)-3B(\lambda)
=
1-4\lambda+10\lambda^2-3\lambda^4
>
1-4\lambda+7\lambda^2>0
\]
on \(0<\lambda<\sqrt{2}-1\). Since \(D(\lambda)>0\) on this interval, this
gives
\[
g(\lambda)<\frac13
\qquad
(0<\lambda<\sqrt{2}-1),
\]
and hence \(m<1/3\).


Let
\[
F(\lambda,\gamma):=\gamma D(\lambda)-B(\lambda).
\]
Since \(F(\lambda,g(\lambda))\equiv0\), we have
$F(\lambda_0,m)=0$. 
Moreover, \(g'(\lambda_0)=0\). Differentiating
\(F(\lambda,g(\lambda))\equiv0\) with respect to \(\lambda\) and then
setting \(\lambda=\lambda_0\), we obtain
$F_\lambda(\lambda_0,m)=0$.

Thus \(\lambda_0\) is a multiple root of \(F(\lambda,m)\), viewed as a
polynomial in \(\lambda\). Since the leading coefficient of \(F(\lambda,m)\)
is \(-3m\ne0\), the quartic discriminant vanishes.


Expanding,
\[
F(\lambda,\gamma)
=
-3\gamma\lambda^4
+(2-6\gamma)\lambda^3
+(4-2\gamma)\lambda^2
+(2\gamma-2)\lambda
+\gamma .
\]
The discriminant with respect to \(\lambda\) is
\[
\operatorname{Disc}_\lambda F(\lambda,\gamma)
=
256Q(\gamma).
\]
Therefore $Q(m)=0$.
Since \(0<m<1/3\) and \(Q\) has a unique root in \((0,1/3)\), we get
$m=t_*$.
Combining this with \eqref{eq:app-transition-gmax}, we have
\begin{equation}\label{eq:app-transition-g-upper}
g(\lambda)\le t_*
\qquad
(0<\lambda<\sqrt{2}-1).
\end{equation}


We shall also use the following feasibility bound. A direct subtraction gives
\begin{equation}\label{eq:app-transition-feasibility-bound}
\frac{2\lambda}{1+\lambda^2}-g(\lambda)
=
\frac{4\lambda^2(2-\lambda-2\lambda^2-\lambda^3)}
{(1+\lambda)^2(1+\lambda^2)(1-3\lambda^2)}
>0
\end{equation}
for \(0<\lambda<\sqrt{2}-1\). Indeed, the denominator is positive because
\(\sqrt{2}-1<1/\sqrt{3}\). Also \(\sqrt{2}-1<1/2\), so
\[
2-\lambda-2\lambda^2-\lambda^3
>
2-\frac12-2\cdot\frac14-\frac18>0.
\]


Now assume first that \(0<\tau<t_*\). Choose \(\lambda=\lambda_0\). Since
\[
\tau<t_*=g(\lambda_0)
<
\frac{2\lambda_0}{1+\lambda_0^2}
\]
by \eqref{eq:app-transition-feasibility-bound}, the value \(\lambda_0\)
satisfies the feasibility condition \eqref{eq:app-transition-feasible}.
Moreover,
\[
T_\tau(\lambda_0)
=
D(\lambda_0)(t_*-\tau)>0,
\]
because \(D(\lambda_0)>0\). Since \(0<\lambda_0<1\), the prefactor in
\eqref{eq:app-transition-compare} is strictly positive. Therefore
\[
\Phi\left(\frac{\tau(1-\lambda_0^2)}{2\lambda_0},\tau\lambda_0\right)
>
c(\tau^2).
\]
Hence
\[
s(\tau^2)>c(\tau^2)
\qquad
(0<\tau<t_*).
\]


Finally assume that \(t_*\le\tau\le1/2\). We show that every feasible
\(\lambda\) gives value at most \(c(\tau^2)\). By
\eqref{eq:app-transition-compare}, it is enough to prove
\[
T_\tau(\lambda)\le0
\qquad
(0<\lambda\le1).
\]


First let \(0<\lambda<\sqrt{2}-1\). Then
\eqref{eq:app-transition-g-upper} gives
\[
g(\lambda)\le t_*\le\tau.
\]
Since \(D(\lambda)>0\), we obtain
\[
T_\tau(\lambda)
=
D(\lambda)(g(\lambda)-\tau)\le0.
\]


Next let \(\sqrt{2}-1\le\lambda<1/\sqrt{3}\). Then
\[
B(\lambda)\le0,
\qquad
D(\lambda)>0,
\]
and therefore
\[
T_\tau(\lambda)=B(\lambda)-\tau D(\lambda)<0.
\]


Finally let \(1/\sqrt{3}\le\lambda\le1\). Then \(D(\lambda)\le0\), so
\(T_\tau(\lambda)=B(\lambda)-\tau D(\lambda)\) is increasing as a function of
\(\tau\). Since \(\tau\le1/2\),
\[
T_\tau(\lambda)\le T_{1/2}(\lambda).
\]
A direct simplification gives
\[
T_{1/2}(\lambda)
=
\frac{\lambda-1}{2}
\left(\lambda(3\lambda^2-1)+5\lambda^2+1\right).
\]
For \(1/\sqrt{3}\le\lambda\le1\), the factor in parentheses is positive, and
\(\lambda-1\le0\). Hence
\[
T_{1/2}(\lambda)\le0,
\]
and so \(T_\tau(\lambda)\le0\).


Thus every feasible point has value at most \(c(\tau^2)\). The endpoint
\(\lambda=1\), equivalently \((x,y)=(0,\tau)\), is feasible and gives
equality. Therefore
\[
s(\tau^2)=c(\tau^2)
\qquad
(t_*\le\tau\le1/2).
\]


Since \(c(\tau^2)=\tau^3(1-\tau)=\beta^{3/2}(1-\sqrt{\beta})\) with
\(\beta=\tau^2\), the desired comparison follows.
\end{proof}

\end{document}
